
Cross-model correlations between alignment benchmarks are routinely interpreted as evidence of alignment-specific coherence, but this interpretation conflates capability-driven co-movement with genuine alignment structure. We demonstrate that controlling for MMLU --- a proxy for general language model capability --- reduces the Spearman correlation between TruthfulQA MC2 (factuality) and a bias proxy across N=296 open-weight LLMs from 0.732 to 0.343, a 53.1\% reduction that is statistically decisive (Fisher z = 6.97, p = 3.22 $\times 10^{-12})$. MMLU alone explains 49.3\% of TruthfulQA variance ($R^2$ = 0.493, p = 3.15 $\times 10^{-45})$ and 76.3\% of bias-proxy variance ($R^2$ = 0.763, p = 5.01 $\times 10^{-94})$, confirming scale confounding of substantial magnitude. A significant residual partial correlation (0.343, BCa 95\% CI [0.180, 0.492], p = 1.40 $\times 10^{-9})$ indicates non-zero alignment-specific co-movement beyond scale, though the structural interpretation remains undetermined at N=296, as the confidence interval spans a pre-specified scenario threshold of 0.40. We apply the Fisher z raw-vs-partial difference test as a diagnostic for scale confounding in alignment benchmark correlation analysis, and demonstrate it through a four-stage pre-registered verification pipeline (H-E1, H-M1, H-M2, H-M3) with 57/57 test cases passing. The bias benchmark used throughout is ARC Challenge accuracy as a proxy for HELM Lite BBQ per-model scores, which were unavailable in the execution environment; all claims regarding factuality-bias coupling are qualified accordingly. An optional RLHF sign test (Tier 3, N=300 base/chat pairs) yields a null result (binomial p = 0.686), providing no evidence that RLHF fine-tuning systematically increases bias-proxy scores.

